\chapter{Podstawy wymuszania i modele generyczne}\label{ch:forcing-foundations}
\chaptermark{Podstawy wymuszania}
\section{Motywacja i założenia metateoretyczne}
W Rozdziale~\ref{ch:constructible} zbudowaliśmy uniwersum
konstruowalne $L$, ograniczając zbiory do tych, które dają się
zdefiniować na kolejnych etapach hierarchii. Wymuszanie służy
odwrotnemu zadaniu: chcemy rozszerzyć model teorii mnogości
o nowy obiekt, zachowując możliwie wiele aksjomatów.
Najprostszy przykład polega na dodaniu ciągu binarnego
przez stopniowe ustalanie jego skończonych fragmentów.

\begin{uwaga}[Model wyjściowy]\label{rem:forcing-background}
W rozdziale pracujemy w zewnętrznym uniwersum ZFC i zakładamy,
że istnieje przeliczalny tranzytywny model zbiorowy
$M\models\mathrm{ZFC}$. Przeliczalność jest własnością
zewnętrzną, nie orzeczeniem modelu $M$ o sobie samym.
Tranzytywność oznacza
$\forall x\in M\ \forall y\in x\ (y\in M)$.
Istnienie takiego modelu nie wynika z samej syntaktycznej
niesprzeczności ZFC. W szczególności poniższych konstrukcji
nie wolno bez dalszego argumentu traktować jako dowodu
względnej niesprzeczności ZFC i niezależności CH.
\end{uwaga}

\section{Warunki i zgodność informacji}
\begin{definicja}[Pojęcie wymuszania]\label{def:forcing-poset}
\emph{Pojęciem wymuszania} nazywamy niepusty zbiór
$\mathbb P$ z relacją $\leq$, która jest zwrotna
i przechodnia. Napis $q\leq p$ oznacza, że $q$ jest
\emph{silniejszym} warunkiem od $p$.
Nie wymagamy antysymetrii. Dla uproszczenia dalszej
konstrukcji zakładamy istnienie największego warunku
$1_{\mathbb P}$, dla którego
$p\leq1_{\mathbb P}$ dla każdego $p\in\mathbb P$.
\end{definicja}

\begin{definicja}[Zgodność i antyłańcuch]\label{def:forcing-compatible}
Warunki $p,q\in\mathbb P$ są zgodne, co zapisujemy
$p\parallel q$, jeżeli istnieje $r\in\mathbb P$
takie, że $r\leq p$ i $r\leq q$.
W przeciwnym przypadku piszemy $p\perp q$.
Zbiór $A\subseteq\mathbb P$ jest antyłańcuchem,
gdy jego dwa różne elementy są niezgodne.
\end{definicja}

\begin{przyklad}[Skończone przybliżenia ciągu binarnego]
\label{ex:cohen-poset}
Niech $\operatorname{Fn}(\omega,2)$ oznacza zbiór
wszystkich funkcji częściowych
$p:F\to\{0,1\}$ o skończonej dziedzinie
$F\subseteq\omega$. Definiujemy
\[
 \mathbb P_{\mathrm C}:=\operatorname{Fn}(\omega,2),
 \qquad q\leq p\iff q\supseteq p.
\]
Największym elementem jest funkcja pusta.
Warunki $p,q$ są zgodne wtedy i tylko wtedy,
gdy zgadzają się na części wspólnej dziedzin:
w takim przypadku ich suma jako funkcji
$p\cup q$ jest wspólnym wzmocnieniem.
Na przykład
$\{(0,0),(3,1)\}$ i $\{(1,1),(3,1)\}$
są zgodne, natomiast
$\{(3,0)\}\perp\{(3,1)\}$.
\end{przyklad}

\begin{figure}[htbp]
\centering
\begin{tikzpicture}[
 level distance=14mm,
 sibling distance=39mm,
 every node/.style={draw,rounded corners,inner sep=3pt,font=\small},
 edge from parent/.style={draw,-{Stealth[length=1.8mm]}}
]
\node {$\varnothing$}
 child {node {$\{(0,0)\}$}
   child {node {$\{(0,0),(1,0)\}$}}
   child {node {$\{(0,0),(1,1)\}$}}}
 child {node {$\{(0,1)\}$}
   child {node {$\{(0,1),(1,0)\}$}}
   child {node {$\{(0,1),(1,1)\}$}}};
\end{tikzpicture}
\caption{Sześć przykładowych wzmocnień warunków Cohena.
Kierunek w dół oznacza silniejsze warunki.}
\label{fig:forcing-cohen-tree}
\end{figure}

\begin{lemat}[Monotoniczność niezgodności]
\label{lem:forcing-compatibility}
Dla dowolnych $p,q,r\in\mathbb P$:
jeśli $p\perp q$ i $r\leq p$, to $r\perp q$.
Jeżeli $p\parallel q$ i $p\leq p'$, $q\leq q'$,
to $p'\parallel q'$.
\end{lemat}
\begin{proof}
Gdyby $r\parallel q$, wspólne wzmocnienie
$s\leq r,q$ spełniałoby też $s\leq p,q$,
wbrew niezgodności. Jeśli natomiast $s\leq p,q$
oraz $p\leq p'$ i $q\leq q'$, to z przechodniości
$s\leq p',q'$, co dowodzi drugiego stwierdzenia.
\end{proof}

\section{Filtry i zbiory gęste}
\begin{definicja}[Filtr]\label{def:forcing-filter}
Filtr w pojęciu wymuszania $(\mathbb P,\leq)$ to niepusty
zbiór $G\subseteq\mathbb P$, który jest domknięty
ku słabszym warunkom i skierowany ku silniejszym:
\begin{enumerate}
\item jeśli $p\in G$ i $p\leq q$, to $q\in G$;
\item jeśli $p,q\in G$, istnieje $r\in G$ z $r\leq p,q$.
\end{enumerate}
\end{definicja}

\begin{definicja}[Gęstość]\label{def:forcing-dense}
Zbiór $D\subseteq\mathbb P$ jest \emph{gęsty},
jeśli $\forall p\in\mathbb P\;\exists q\in D\ (q\leq p)$.
Nazywamy go \emph{otwartym w porządku},
jeżeli $p\in D$ i $q\leq p$ implikują $q\in D$.
\end{definicja}

\begin{lemat}[Otwarte domknięcie zbioru gęstego]
\label{lem:forcing-open-dense}
Jeżeli $D$ jest gęsty, to
$D^\downarrow=\{p:\exists q\in D\ (p\leq q)\}$
jest gęsty i otwarty w porządku. Dla dowolnego
filtru $G$ zachodzi
$G\cap D\ne\varnothing$ wtedy i tylko wtedy,
gdy $G\cap D^\downarrow\ne\varnothing$.
\end{lemat}
\begin{proof}
Jeśli $p\in\mathbb P$, gęstość daje
$q\in D$ z $q\leq p$. Zatem $q\in D^\downarrow$,
więc jest gęsty. Gdy $r\leq p\in D^\downarrow$
oraz $p\leq q\in D$, to $r\leq q$, więc
$r\in D^\downarrow$. Jedna implikacja dotycząca
filtrów wynika z $D\subseteq D^\downarrow$.
Dla drugiej, z $p\in G\cap D^\downarrow$
wybieramy $q\in D$ z $p\leq q$; wtedy
$q\in G\cap D$, ponieważ filtr jest
domknięty ku słabszym warunkom.
\end{proof}

\begin{definicja}[Filtr generyczny]\label{def:forcing-generic}
Niech $M$ będzie tranzytywnym modelem ZFC,
$(\mathbb P,\leq)\in M$ pojęciem wymuszania.
Filtr $G\subseteq\mathbb P$ nazywamy
\emph{$M$-generycznym}, jeśli dla każdego
gęstego $D\subseteq\mathbb P$ należącego do
$M$ zachodzi $G\cap D\ne\varnothing$.
Nie wymagamy $G\in M$. Gęstość $D\in M$ jest
absolutna, gdyż wszystkie kwantyfikatory
biegną po $\mathbb P\subseteq M$.
\end{definicja}

\begin{lemat}[Generyczność rozstrzyga zgodność]
\label{lem:forcing-generic-decision}
Jeśli $G$ jest filtrem $M$-generycznym, to
dla każdego $p\in\mathbb P$ albo $p\in G$,
albo istnieje $q\in G$ niezgodny z $p$.
Dokładnie jedna z tych możliwości zachodzi.
\end{lemat}
\begin{proof}
Z tranzytywności $p\in M$. Zbiór
$D_p=\{q\in\mathbb P:q\leq p\text{ lub }q\perp p\}$
należy do $M$ i jest gęsty. Dla każdego $r$
albo $r\perp p$, albo istnieje wspólne
wzmocnienie $s\leq r,p$. Generyczność
daje $q\in G\cap D_p$. Jeśli $q\leq p$,
to $p\in G$ przez domknięcie ku słabszym
warunkom. W drugim przypadku $q\perp p$.
Oba przypadki nie mogą zachodzić razem,
gdyż dwa elementy filtru są zgodne.
\end{proof}

\section{Lemat Rasiowej--Sikorskiego}
\begin{twierdzenie}[Lemat Rasiowej--Sikorskiego]
\label{thm:rasiowa-sikorski}
Niech $\mathbb P$ będzie niepustym
przeliczalnym pojęciem wymuszania,
$(D_n)_{n\in\omega}$ ciągiem gęstych
podzbiorów $\mathbb P$, a $p_*\in\mathbb P$.
W ZFC istnieje filtr $G\subseteq\mathbb P$
zawierający $p_*$ i przecinający
każdy zbiór $D_n$.
\end{twierdzenie}
\begin{proof}
Ustalmy enumerację $\mathbb P$. Definiujemy
$p_0=p_*$, a $p_{n+1}$ jako pierwszy
w tej enumeracji element $D_n$ spełniający
$p_{n+1}\leq p_n$. Gęstość zapewnia istnienie.
Niech
\[
 G=\{q\in\mathbb P:\exists n\in\omega\ (p_n\leq q)\}.
\]
Zbiór zawiera $p_0$, więc jest niepusty.
Jeżeli $q\in G$ i $q\leq r$,
to $p_n\leq q\leq r$ dla pewnego $n$,
czyli $r\in G$. Jeśli $q,r\in G$
i $p_n\leq q$, $p_m\leq r$, to
$p_{\max(m,n)}\in G$ jest wspólnym
wzmocnieniem obu warunków. Zatem $G$
jest filtrem. Ponieważ $p_{n+1}\in G\cap D_n$,
przecina wszystkie wskazane zbiory.
\end{proof}

\begin{wniosek}[Filtr generyczny nad modelem przeliczalnym]
\label{cor:countable-generic}
Jeżeli $M$ jest przeliczalnym tranzytywnym
modelem ZFC, $(\mathbb P,\leq)\in M$
i $p_*\in\mathbb P$, to istnieje zewnętrzny
filtr $M$-generyczny zawierający $p_*$.
\end{wniosek}
\begin{proof}
Tranzytywność zapewnia $\mathbb P\subseteq M$,
zatem $\mathbb P$ jest przeliczalny zewnętrznie.
Rodzina gęstych podzbiorów $\mathbb P$
należących do $M$ jest także przeliczalna,
gdyż jest podzbiorem $M$.
Wyliczmy ją jako $(D_n)$, powtarzając
elementy w przypadku rodziny skończonej.
Twierdzenie~\ref{thm:rasiowa-sikorski}
daje odpowiedni filtr.
\end{proof}

\begin{uwaga}[Zewnętrzna perspektywa]
\label{rem:forcing-external-generic}
Wyliczenie gęstych zbiorów należących do $M$
jest dokonane na zewnątrz i nie musi należeć
do samego $M$. Nie dowodzi więc istnienia
generycznego filtru jako elementu $M$.
Nie wymaga też przecięcia wszystkich
gęstych zbiorów zewnętrznego uniwersum.
\end{uwaga}
